\documentclass[12pt]{article}
\usepackage[utf8]{inputenc}
\usepackage[T1]{fontenc}
\usepackage{hyperref}
\title{Urat Language Resources for K-12 Curriculum Development}
\author{Compiled by Claude Sonnet 5 (claude-sonnet-5) for Andr\'as Kornai}
\date{}
\begin{document}
\maketitle

\section{Introduction}
Urat (ISO 639-3 urt) is a Torricelli language of the Urim subgroup, spoken in the Dreikikir district of East Sepik Province, Papua New Guinea. Sources give roughly 7,000 speakers, and Wikipedia lists four dialects (North, Central, South, East). The dictionary and New Testament found are based on the Yehre dialect.

\section{Primary Educational Resources}
No primers, textbooks or school programmes were found in search results.

\section{Linguistic and Reference Materials}
\begin{itemize}
\item Urat--Melanesian Pidgin--English dictionary (Webonary, SIL PNG, Yehre dialect, not exhaustive): \url{https://www.webonary.work/urat/overview/introduction/}
\item Copyright and licence page for the dictionary: \url{https://www.webonary.work/urat/overview/copyright/}
\item Wikipedia entry: \url{https://en.wikipedia.org/wiki/Urat_language}
\item Urim languages overview: \url{https://en.wikipedia.org/wiki/Urim_languages}
\end{itemize}

\section{Religious and Community Resources}
The Urat New Testament, \emph{Wusyep Bwore tikin Got} (2016, SIL PNG and Wycliffe), is available from eBible.org under CC BY-ND 4.0:
\begin{itemize}
\item \url{https://ebible.org/pdf/urt/urt_all.pdf}
\item \url{https://ebible.org/epub/urt.epub}
\item \url{https://dbs.org/bibles/URTWBT}
\end{itemize}

\section{Digital and Audio Resources}
The eBible.org and Webonary resources above are digital. No audio or video was found.

\section{Curriculum Development Context}
Urat has a dictionary and a scripture translation in a written form, which could serve as a base for an orthography and early-literacy materials. The ND licence on the scripture text restricts adapting it. Coverage of the non-Yehre dialects is unclear.
\end{document}